% =============================================================================
%   Stability Monitoring Methods Section
% =============================================================================
%   This file is a fragment, not a standalone document. Drop it into your
%   paper with % =============================================================================
%   Stability Monitoring Methods Section
% =============================================================================
%   This file is a fragment, not a standalone document. Drop it into your
%   paper with % =============================================================================
%   Stability Monitoring Methods Section
% =============================================================================
%   This file is a fragment, not a standalone document. Drop it into your
%   paper with % =============================================================================
%   Stability Monitoring Methods Section
% =============================================================================
%   This file is a fragment, not a standalone document. Drop it into your
%   paper with \input{stability_monitoring_methods.tex}, or paste directly.
%
%   Required packages in your main preamble:
%       \usepackage{amsmath, amssymb, amsfonts}
%       \usepackage{bbm}    % for \mathbbm{1} indicator function
%       \usepackage{bm}     % for \bm{} bold math
%
%   If \mathbbm is unavailable in your TeX distribution, substitute:
%       \usepackage{dsfont}     and use \mathds{1}
%   or fall back to \mathbf{1} (no extra package).
% =============================================================================

\section{Methods: Online Stability Monitoring}

We monitor the rollouts of a vision--language--action (VLA) student policy on
the 6-DOF SO-101 arm and intervene when classical control-theoretic
indicators of trajectory instability exceed thresholds calibrated on a teacher
policy's successful rollouts. All indicators are derived solely from the data
that the LeRobot \texttt{record} pipeline persists by default, so the monitor
imposes no additional logging requirements and runs at a small fraction of the
cost of a second VLA inference pass.

\subsection{Notation and Data}

The SO-101 has six Feetech STS3215 serial-bus servos, indexed
$i \in \{1,\ldots,6\}$ in the order
\textit{shoulder\_pan}, \textit{shoulder\_lift}, \textit{elbow\_flex},
\textit{wrist\_flex}, \textit{wrist\_roll}, \textit{gripper}.
The default \texttt{lerobot-record} schema persists, at sampling rate
$f_s = 30$~Hz with period $T_s = 1/f_s$, two synchronized vectors per frame:
\begin{align}
    \mathbf{a}[k] &= \big[a_1[k], \ldots, a_6[k]\big]^\top \in \mathbb{R}^6 && \text{(commanded joint positions; \texttt{action})} \\
    \mathbf{q}[k] &= \big[q_1[k], \ldots, q_6[k]\big]^\top \in \mathbb{R}^6 && \text{(measured joint positions; \texttt{observation.state})}
\end{align}
The motors expose load, current, velocity, and temperature feedback over the
serial bus, but these channels are not logged by the default recorder; the
metrics below therefore use $\mathbf{a}$ and $\mathbf{q}$ exclusively.

For real-time monitoring we compute all metrics over a sliding window of
length $N$ samples ending at the current step $k$,
\begin{equation}
    \mathcal{W}[k] \;=\; \{k - N + 1, \ldots, k\}, \qquad \Delta \;=\; N T_s.
\end{equation}
The 12-bit magnetic encoder on the STS3215 quantises position at $\approx
0.088^\circ$ and the gearbox introduces up to $0.5^\circ$ of backlash, so
naive numerical differentiation of $\mathbf{q}$ is dominated by quantisation
noise. We therefore apply a Savitzky--Golay smoother (window 7, polynomial
order 3) to $\mathbf{q}[k]$ before differentiation, denoting the smoothed
signal $\tilde{\mathbf{q}}[k]$.

\subsection{Position Tracking Error}

\textbf{Motivation.} Because the SO-101 servos accept position setpoints
rather than torques, the gap between commanded and measured position is the
direct, controller-level analogue of a torque error. It captures both
transient failures (the policy commanding more motion than the joint can
deliver under load) and persistent stalls (a joint pinned at its load limit).
Since both signals are recorded synchronously and at the same rate, this is
the cleanest single stability indicator available without extending the
recorder.

Define the per-joint tracking error with a one-step offset to align the
action with the resulting state:
\begin{equation}
    e_i[k] \;=\; a_i[k-1] - q_i[k].
\end{equation}
Per-joint windowed root-mean-square error:
\begin{equation}
    E_i^{\mathrm{RMS}}[k] \;=\; \sqrt{\,\frac{1}{N} \sum_{j \in \mathcal{W}[k]} e_i[j]^2\,}.
\end{equation}
Aggregate norm-error across joints:
\begin{equation}
    E^{\mathrm{RMS}}[k] \;=\; \sqrt{\,\sum_{i=1}^{6} \big(E_i^{\mathrm{RMS}}[k]\big)^2\,}.
\end{equation}

\subsection{Jerk-Based Smoothness}

\textbf{Motivation.} Jerk---the third time-derivative of position---is the
canonical scalar quantification of trajectory roughness, with a long history
in human motor-control literature and in robotics. A policy beginning to lose
stability typically produces jerky motion well before it produces visibly
incorrect motion, providing the earliest-firing indicator in our suite. We
compute jerk on the smoothed measured signal to assess executed motion;
computing it independently on $\mathbf{a}[k]$ (for which no smoothing is
needed) and comparing the two streams further distinguishes
controller-tracking roughness from policy-output roughness.

Discrete third-derivative approximation via a five-point central
finite-difference stencil on the smoothed signal:
\begin{equation}
    \dddot{\tilde{q}}_i[k]
    \;\approx\;
    \frac{\tilde{q}_i[k+2] - 2\,\tilde{q}_i[k+1] + 2\,\tilde{q}_i[k-1] - \tilde{q}_i[k-2]}{2\,T_s^3}.
\end{equation}
Per-joint windowed RMS jerk and its aggregate:
\begin{equation}
    J_i^{\mathrm{RMS}}[k] \;=\; \sqrt{\,\frac{1}{N} \sum_{j \in \mathcal{W}[k]} \dddot{\tilde{q}}_i[j]^2\,},
    \qquad
    J^{\mathrm{RMS}}[k] \;=\; \sqrt{\,\sum_{i=1}^{6} \big(J_i^{\mathrm{RMS}}[k]\big)^2\,}.
\end{equation}

\subsection{Spectral Arc Length (SPARC)}

\textbf{Motivation.} RMS jerk is sensitive to movement scale and duration,
which makes cross-trajectory comparison brittle. SPARC, the negated arc
length of the normalised speed spectrum, is dimensionless and approximately
invariant to amplitude and time-scale changes. It is the most reliable
single-trajectory smoothness statistic in the motor-control literature and
provides a complementary view to RMS jerk: trajectories that look smooth in
RMS but contain narrow-band oscillations are penalised by SPARC.

Define the speed profile
\begin{equation}
    v[k] \;=\; \big\|\, \dot{\tilde{\mathbf{q}}}[k] \,\big\|_2,
\end{equation}
and let $V(\omega)$ denote the discrete-time Fourier transform of $v$ over
the window $\mathcal{W}[k]$, with $\hat{V}(\omega) = V(\omega)/V(0)$ its
normalised magnitude. SPARC is then
\begin{equation}
    \eta_{\mathrm{SPARC}}[k]
    \;=\;
    -\int_{0}^{\omega_c}
    \sqrt{\left(\frac{1}{\omega_c}\right)^{\!2} \!+\! \left(\frac{d\hat{V}(\omega)}{d\omega}\right)^{\!2}}
    \, d\omega,
\end{equation}
with the adaptive cutoff
\begin{equation}
    \omega_c \;=\; \min\!\Big\{\, \omega_c^{\max}, \; \min\!\big\{\omega \,\big|\, \hat{V}(r) < \bar{V}\ \forall\, r > \omega \big\} \,\Big\}.
\end{equation}
At $f_s = 30$~Hz the Nyquist frequency is $15$~Hz, so we use
$\omega_c^{\max} = 12$~Hz and $\bar{V} = 0.05$. Smoother trajectories yield
$\eta_{\mathrm{SPARC}}$ closer to zero from below; rougher motion drives it
more negative.

\subsection{High-Frequency Power Ratio of the Tracking Error}

\textbf{Motivation.} When a position controller cannot follow its setpoint,
the error signal develops oscillatory content well above the bandwidth of the
intended motion---the spectral fingerprint of chatter, limit-cycling, and
incipient instability. On a torque-controlled arm this is conventionally
detected on the commanded torque; lacking torque telemetry, we apply the same
analysis to $e_i[k]$, which carries the same information because closed-loop
oscillation manifests as periodic error around the setpoint.

Let $S_{e_i}(\omega, k)$ denote the Welch power spectral density estimate of
$e_i$ computed on $\mathcal{W}[k]$, with crossover frequency $\omega_h$:
\begin{equation}
    \rho_i^{\mathrm{HF}}[k]
    \;=\;
    \frac{\displaystyle \int_{\omega_h}^{\omega_{\mathrm{Nyq}}} S_{e_i}(\omega, k)\, d\omega}{\displaystyle \int_{0}^{\omega_{\mathrm{Nyq}}} S_{e_i}(\omega, k)\, d\omega}.
\end{equation}
The aggregate is the worst-joint ratio,
\begin{equation}
    \rho^{\mathrm{HF}}[k] \;=\; \max_{i \in \{1,\ldots,6\}} \rho_i^{\mathrm{HF}}[k].
\end{equation}
For SO-101 we use $\omega_h = 5$~Hz and $\omega_{\mathrm{Nyq}} = 15$~Hz.
Sustained $\rho^{\mathrm{HF}} > 0.2$ is a strong chatter indicator on
arm-scale dynamics.

\subsection{Stall Indicator and Joint-Limit Clipping}

\textbf{Motivation.} On a torque-instrumented arm we would threshold the
commanded torque against the motor's torque limit. Without that telemetry we
detect saturation behaviourally: when the policy commands meaningful motion
but the joint fails to move, the motor is effectively at its load limit. We
combine this with a separate detector for action saturation against the
calibrated joint range.

Let $\delta_a^{(i)}, \delta_q^{(i)}$ be per-joint command and motion
thresholds, calibrated as the 95th percentile of the corresponding magnitudes
on a corpus of successful teacher rollouts. Over a short look-back of $M$
samples, the instantaneous stall indicator is
\begin{equation}
    \sigma_i^{\mathrm{stall}}[k] \;=\; \mathbbm{1}\!\left[\; |a_i[k] - a_i[k-M]| > \delta_a^{(i)} \;\wedge\; |q_i[k] - q_i[k-M]| < \delta_q^{(i)} \;\right].
\end{equation}
With calibrated joint range $[a_i^{\min}, a_i^{\max}]$ and tolerance
$\epsilon_i = 0.02 \cdot (a_i^{\max} - a_i^{\min})$, the clipping indicator is
\begin{equation}
    \sigma_i^{\mathrm{clip}}[k] \;=\; \mathbbm{1}\!\left[\; a_i[k] \le a_i^{\min} + \epsilon_i \;\vee\; a_i[k] \ge a_i^{\max} - \epsilon_i \;\right].
\end{equation}
The combined per-joint flag and its windowed aggregate are
\begin{align}
    \sigma_i[k] &\;=\; \sigma_i^{\mathrm{stall}}[k] \;\vee\; \sigma_i^{\mathrm{clip}}[k], \\
    \bar{\sigma}[k] &\;=\; \frac{1}{N} \sum_{j \in \mathcal{W}[k]} \mathbbm{1}\!\left[\, \max_{i} \sigma_i[j] = 1 \,\right].
\end{align}
We use $M = 5$ ($\approx 167$~ms at 30~Hz), short enough to react to abrupt
stalls but long enough to ignore single-frame quantisation jitter.

\subsection{Combined Intervention Rule}

We assemble the five sign-aligned signals (larger means worse) into a single
monitor vector,
\begin{equation}
    \mathbf{m}[k] \;=\; \big[\, E^{\mathrm{RMS}}[k], \; J^{\mathrm{RMS}}[k], \; -\eta_{\mathrm{SPARC}}[k], \; \rho^{\mathrm{HF}}[k], \; \bar{\sigma}[k] \,\big]^\top \in \mathbb{R}^5.
\end{equation}
A threshold vector $\bm{\theta} \in \mathbb{R}^5$ is calibrated as the 95th
percentile of each component over a held-out corpus of successful teacher
rollouts. The any-of intervention rule fires when any single channel
exceeds its threshold,
\begin{equation}
    \mathcal{I}^{\mathrm{any}}[k]
    \;=\;
    \mathbbm{1}\!\left[\, \exists\, \ell \in \{1,\ldots,5\} : m_\ell[k] > \theta_\ell \,\right],
\end{equation}
while the weighted-sum rule, with non-negative weights
$\mathbf{w} \in \mathbb{R}_{\ge 0}^{5}$ and aggregate threshold
$\theta_{\Sigma}$, integrates evidence across channels:
\begin{equation}
    \mathcal{I}^{\Sigma}[k]
    \;=\;
    \mathbbm{1}\!\left[\, \sum_{\ell=1}^{5} w_\ell \, \frac{m_\ell[k]}{\theta_\ell} > \theta_{\Sigma} \,\right].
\end{equation}
The any-of rule fires earlier and is easier to attribute when debugging; the
weighted-sum rule is more robust to single-channel noise. To suppress
single-frame artefacts, we additionally require the chosen rule to remain
active for $P$ consecutive samples before triggering an intervention,
\begin{equation}
    \mathcal{I}^{\star}[k] \;=\; \prod_{p = 0}^{P-1} \mathcal{I}[k - p],
\end{equation}
with $P = 3$ at 30~Hz ($\approx 100$~ms persistence).

\paragraph{Optional extension.} If \texttt{Present\_Load} or
\texttt{Present\_Current} is added to the recorded observation features by
extending the LeRobot \texttt{Robot} interface, the stall indicator
$\sigma_i^{\mathrm{stall}}$ can be replaced by the conventional torque
saturation rate, and the high-frequency power ratio can be recomputed on the
load signal directly, recovering the full classical signal set without
modifying the rest of the pipeline.
, or paste directly.
%
%   Required packages in your main preamble:
%       \usepackage{amsmath, amssymb, amsfonts}
%       \usepackage{bbm}    % for \mathbbm{1} indicator function
%       \usepackage{bm}     % for \bm{} bold math
%
%   If \mathbbm is unavailable in your TeX distribution, substitute:
%       \usepackage{dsfont}     and use \mathds{1}
%   or fall back to \mathbf{1} (no extra package).
% =============================================================================

\section{Methods: Online Stability Monitoring}

We monitor the rollouts of a vision--language--action (VLA) student policy on
the 6-DOF SO-101 arm and intervene when classical control-theoretic
indicators of trajectory instability exceed thresholds calibrated on a teacher
policy's successful rollouts. All indicators are derived solely from the data
that the LeRobot \texttt{record} pipeline persists by default, so the monitor
imposes no additional logging requirements and runs at a small fraction of the
cost of a second VLA inference pass.

\subsection{Notation and Data}

The SO-101 has six Feetech STS3215 serial-bus servos, indexed
$i \in \{1,\ldots,6\}$ in the order
\textit{shoulder\_pan}, \textit{shoulder\_lift}, \textit{elbow\_flex},
\textit{wrist\_flex}, \textit{wrist\_roll}, \textit{gripper}.
The default \texttt{lerobot-record} schema persists, at sampling rate
$f_s = 30$~Hz with period $T_s = 1/f_s$, two synchronized vectors per frame:
\begin{align}
    \mathbf{a}[k] &= \big[a_1[k], \ldots, a_6[k]\big]^\top \in \mathbb{R}^6 && \text{(commanded joint positions; \texttt{action})} \\
    \mathbf{q}[k] &= \big[q_1[k], \ldots, q_6[k]\big]^\top \in \mathbb{R}^6 && \text{(measured joint positions; \texttt{observation.state})}
\end{align}
The motors expose load, current, velocity, and temperature feedback over the
serial bus, but these channels are not logged by the default recorder; the
metrics below therefore use $\mathbf{a}$ and $\mathbf{q}$ exclusively.

For real-time monitoring we compute all metrics over a sliding window of
length $N$ samples ending at the current step $k$,
\begin{equation}
    \mathcal{W}[k] \;=\; \{k - N + 1, \ldots, k\}, \qquad \Delta \;=\; N T_s.
\end{equation}
The 12-bit magnetic encoder on the STS3215 quantises position at $\approx
0.088^\circ$ and the gearbox introduces up to $0.5^\circ$ of backlash, so
naive numerical differentiation of $\mathbf{q}$ is dominated by quantisation
noise. We therefore apply a Savitzky--Golay smoother (window 7, polynomial
order 3) to $\mathbf{q}[k]$ before differentiation, denoting the smoothed
signal $\tilde{\mathbf{q}}[k]$.

\subsection{Position Tracking Error}

\textbf{Motivation.} Because the SO-101 servos accept position setpoints
rather than torques, the gap between commanded and measured position is the
direct, controller-level analogue of a torque error. It captures both
transient failures (the policy commanding more motion than the joint can
deliver under load) and persistent stalls (a joint pinned at its load limit).
Since both signals are recorded synchronously and at the same rate, this is
the cleanest single stability indicator available without extending the
recorder.

Define the per-joint tracking error with a one-step offset to align the
action with the resulting state:
\begin{equation}
    e_i[k] \;=\; a_i[k-1] - q_i[k].
\end{equation}
Per-joint windowed root-mean-square error:
\begin{equation}
    E_i^{\mathrm{RMS}}[k] \;=\; \sqrt{\,\frac{1}{N} \sum_{j \in \mathcal{W}[k]} e_i[j]^2\,}.
\end{equation}
Aggregate norm-error across joints:
\begin{equation}
    E^{\mathrm{RMS}}[k] \;=\; \sqrt{\,\sum_{i=1}^{6} \big(E_i^{\mathrm{RMS}}[k]\big)^2\,}.
\end{equation}

\subsection{Jerk-Based Smoothness}

\textbf{Motivation.} Jerk---the third time-derivative of position---is the
canonical scalar quantification of trajectory roughness, with a long history
in human motor-control literature and in robotics. A policy beginning to lose
stability typically produces jerky motion well before it produces visibly
incorrect motion, providing the earliest-firing indicator in our suite. We
compute jerk on the smoothed measured signal to assess executed motion;
computing it independently on $\mathbf{a}[k]$ (for which no smoothing is
needed) and comparing the two streams further distinguishes
controller-tracking roughness from policy-output roughness.

Discrete third-derivative approximation via a five-point central
finite-difference stencil on the smoothed signal:
\begin{equation}
    \dddot{\tilde{q}}_i[k]
    \;\approx\;
    \frac{\tilde{q}_i[k+2] - 2\,\tilde{q}_i[k+1] + 2\,\tilde{q}_i[k-1] - \tilde{q}_i[k-2]}{2\,T_s^3}.
\end{equation}
Per-joint windowed RMS jerk and its aggregate:
\begin{equation}
    J_i^{\mathrm{RMS}}[k] \;=\; \sqrt{\,\frac{1}{N} \sum_{j \in \mathcal{W}[k]} \dddot{\tilde{q}}_i[j]^2\,},
    \qquad
    J^{\mathrm{RMS}}[k] \;=\; \sqrt{\,\sum_{i=1}^{6} \big(J_i^{\mathrm{RMS}}[k]\big)^2\,}.
\end{equation}

\subsection{Spectral Arc Length (SPARC)}

\textbf{Motivation.} RMS jerk is sensitive to movement scale and duration,
which makes cross-trajectory comparison brittle. SPARC, the negated arc
length of the normalised speed spectrum, is dimensionless and approximately
invariant to amplitude and time-scale changes. It is the most reliable
single-trajectory smoothness statistic in the motor-control literature and
provides a complementary view to RMS jerk: trajectories that look smooth in
RMS but contain narrow-band oscillations are penalised by SPARC.

Define the speed profile
\begin{equation}
    v[k] \;=\; \big\|\, \dot{\tilde{\mathbf{q}}}[k] \,\big\|_2,
\end{equation}
and let $V(\omega)$ denote the discrete-time Fourier transform of $v$ over
the window $\mathcal{W}[k]$, with $\hat{V}(\omega) = V(\omega)/V(0)$ its
normalised magnitude. SPARC is then
\begin{equation}
    \eta_{\mathrm{SPARC}}[k]
    \;=\;
    -\int_{0}^{\omega_c}
    \sqrt{\left(\frac{1}{\omega_c}\right)^{\!2} \!+\! \left(\frac{d\hat{V}(\omega)}{d\omega}\right)^{\!2}}
    \, d\omega,
\end{equation}
with the adaptive cutoff
\begin{equation}
    \omega_c \;=\; \min\!\Big\{\, \omega_c^{\max}, \; \min\!\big\{\omega \,\big|\, \hat{V}(r) < \bar{V}\ \forall\, r > \omega \big\} \,\Big\}.
\end{equation}
At $f_s = 30$~Hz the Nyquist frequency is $15$~Hz, so we use
$\omega_c^{\max} = 12$~Hz and $\bar{V} = 0.05$. Smoother trajectories yield
$\eta_{\mathrm{SPARC}}$ closer to zero from below; rougher motion drives it
more negative.

\subsection{High-Frequency Power Ratio of the Tracking Error}

\textbf{Motivation.} When a position controller cannot follow its setpoint,
the error signal develops oscillatory content well above the bandwidth of the
intended motion---the spectral fingerprint of chatter, limit-cycling, and
incipient instability. On a torque-controlled arm this is conventionally
detected on the commanded torque; lacking torque telemetry, we apply the same
analysis to $e_i[k]$, which carries the same information because closed-loop
oscillation manifests as periodic error around the setpoint.

Let $S_{e_i}(\omega, k)$ denote the Welch power spectral density estimate of
$e_i$ computed on $\mathcal{W}[k]$, with crossover frequency $\omega_h$:
\begin{equation}
    \rho_i^{\mathrm{HF}}[k]
    \;=\;
    \frac{\displaystyle \int_{\omega_h}^{\omega_{\mathrm{Nyq}}} S_{e_i}(\omega, k)\, d\omega}{\displaystyle \int_{0}^{\omega_{\mathrm{Nyq}}} S_{e_i}(\omega, k)\, d\omega}.
\end{equation}
The aggregate is the worst-joint ratio,
\begin{equation}
    \rho^{\mathrm{HF}}[k] \;=\; \max_{i \in \{1,\ldots,6\}} \rho_i^{\mathrm{HF}}[k].
\end{equation}
For SO-101 we use $\omega_h = 5$~Hz and $\omega_{\mathrm{Nyq}} = 15$~Hz.
Sustained $\rho^{\mathrm{HF}} > 0.2$ is a strong chatter indicator on
arm-scale dynamics.

\subsection{Stall Indicator and Joint-Limit Clipping}

\textbf{Motivation.} On a torque-instrumented arm we would threshold the
commanded torque against the motor's torque limit. Without that telemetry we
detect saturation behaviourally: when the policy commands meaningful motion
but the joint fails to move, the motor is effectively at its load limit. We
combine this with a separate detector for action saturation against the
calibrated joint range.

Let $\delta_a^{(i)}, \delta_q^{(i)}$ be per-joint command and motion
thresholds, calibrated as the 95th percentile of the corresponding magnitudes
on a corpus of successful teacher rollouts. Over a short look-back of $M$
samples, the instantaneous stall indicator is
\begin{equation}
    \sigma_i^{\mathrm{stall}}[k] \;=\; \mathbbm{1}\!\left[\; |a_i[k] - a_i[k-M]| > \delta_a^{(i)} \;\wedge\; |q_i[k] - q_i[k-M]| < \delta_q^{(i)} \;\right].
\end{equation}
With calibrated joint range $[a_i^{\min}, a_i^{\max}]$ and tolerance
$\epsilon_i = 0.02 \cdot (a_i^{\max} - a_i^{\min})$, the clipping indicator is
\begin{equation}
    \sigma_i^{\mathrm{clip}}[k] \;=\; \mathbbm{1}\!\left[\; a_i[k] \le a_i^{\min} + \epsilon_i \;\vee\; a_i[k] \ge a_i^{\max} - \epsilon_i \;\right].
\end{equation}
The combined per-joint flag and its windowed aggregate are
\begin{align}
    \sigma_i[k] &\;=\; \sigma_i^{\mathrm{stall}}[k] \;\vee\; \sigma_i^{\mathrm{clip}}[k], \\
    \bar{\sigma}[k] &\;=\; \frac{1}{N} \sum_{j \in \mathcal{W}[k]} \mathbbm{1}\!\left[\, \max_{i} \sigma_i[j] = 1 \,\right].
\end{align}
We use $M = 5$ ($\approx 167$~ms at 30~Hz), short enough to react to abrupt
stalls but long enough to ignore single-frame quantisation jitter.

\subsection{Combined Intervention Rule}

We assemble the five sign-aligned signals (larger means worse) into a single
monitor vector,
\begin{equation}
    \mathbf{m}[k] \;=\; \big[\, E^{\mathrm{RMS}}[k], \; J^{\mathrm{RMS}}[k], \; -\eta_{\mathrm{SPARC}}[k], \; \rho^{\mathrm{HF}}[k], \; \bar{\sigma}[k] \,\big]^\top \in \mathbb{R}^5.
\end{equation}
A threshold vector $\bm{\theta} \in \mathbb{R}^5$ is calibrated as the 95th
percentile of each component over a held-out corpus of successful teacher
rollouts. The any-of intervention rule fires when any single channel
exceeds its threshold,
\begin{equation}
    \mathcal{I}^{\mathrm{any}}[k]
    \;=\;
    \mathbbm{1}\!\left[\, \exists\, \ell \in \{1,\ldots,5\} : m_\ell[k] > \theta_\ell \,\right],
\end{equation}
while the weighted-sum rule, with non-negative weights
$\mathbf{w} \in \mathbb{R}_{\ge 0}^{5}$ and aggregate threshold
$\theta_{\Sigma}$, integrates evidence across channels:
\begin{equation}
    \mathcal{I}^{\Sigma}[k]
    \;=\;
    \mathbbm{1}\!\left[\, \sum_{\ell=1}^{5} w_\ell \, \frac{m_\ell[k]}{\theta_\ell} > \theta_{\Sigma} \,\right].
\end{equation}
The any-of rule fires earlier and is easier to attribute when debugging; the
weighted-sum rule is more robust to single-channel noise. To suppress
single-frame artefacts, we additionally require the chosen rule to remain
active for $P$ consecutive samples before triggering an intervention,
\begin{equation}
    \mathcal{I}^{\star}[k] \;=\; \prod_{p = 0}^{P-1} \mathcal{I}[k - p],
\end{equation}
with $P = 3$ at 30~Hz ($\approx 100$~ms persistence).

\paragraph{Optional extension.} If \texttt{Present\_Load} or
\texttt{Present\_Current} is added to the recorded observation features by
extending the LeRobot \texttt{Robot} interface, the stall indicator
$\sigma_i^{\mathrm{stall}}$ can be replaced by the conventional torque
saturation rate, and the high-frequency power ratio can be recomputed on the
load signal directly, recovering the full classical signal set without
modifying the rest of the pipeline.
, or paste directly.
%
%   Required packages in your main preamble:
%       \usepackage{amsmath, amssymb, amsfonts}
%       \usepackage{bbm}    % for \mathbbm{1} indicator function
%       \usepackage{bm}     % for \bm{} bold math
%
%   If \mathbbm is unavailable in your TeX distribution, substitute:
%       \usepackage{dsfont}     and use \mathds{1}
%   or fall back to \mathbf{1} (no extra package).
% =============================================================================

\section{Methods: Online Stability Monitoring}

We monitor the rollouts of a vision--language--action (VLA) student policy on
the 6-DOF SO-101 arm and intervene when classical control-theoretic
indicators of trajectory instability exceed thresholds calibrated on a teacher
policy's successful rollouts. All indicators are derived solely from the data
that the LeRobot \texttt{record} pipeline persists by default, so the monitor
imposes no additional logging requirements and runs at a small fraction of the
cost of a second VLA inference pass.

\subsection{Notation and Data}

The SO-101 has six Feetech STS3215 serial-bus servos, indexed
$i \in \{1,\ldots,6\}$ in the order
\textit{shoulder\_pan}, \textit{shoulder\_lift}, \textit{elbow\_flex},
\textit{wrist\_flex}, \textit{wrist\_roll}, \textit{gripper}.
The default \texttt{lerobot-record} schema persists, at sampling rate
$f_s = 30$~Hz with period $T_s = 1/f_s$, two synchronized vectors per frame:
\begin{align}
    \mathbf{a}[k] &= \big[a_1[k], \ldots, a_6[k]\big]^\top \in \mathbb{R}^6 && \text{(commanded joint positions; \texttt{action})} \\
    \mathbf{q}[k] &= \big[q_1[k], \ldots, q_6[k]\big]^\top \in \mathbb{R}^6 && \text{(measured joint positions; \texttt{observation.state})}
\end{align}
The motors expose load, current, velocity, and temperature feedback over the
serial bus, but these channels are not logged by the default recorder; the
metrics below therefore use $\mathbf{a}$ and $\mathbf{q}$ exclusively.

For real-time monitoring we compute all metrics over a sliding window of
length $N$ samples ending at the current step $k$,
\begin{equation}
    \mathcal{W}[k] \;=\; \{k - N + 1, \ldots, k\}, \qquad \Delta \;=\; N T_s.
\end{equation}
The 12-bit magnetic encoder on the STS3215 quantises position at $\approx
0.088^\circ$ and the gearbox introduces up to $0.5^\circ$ of backlash, so
naive numerical differentiation of $\mathbf{q}$ is dominated by quantisation
noise. We therefore apply a Savitzky--Golay smoother (window 7, polynomial
order 3) to $\mathbf{q}[k]$ before differentiation, denoting the smoothed
signal $\tilde{\mathbf{q}}[k]$.

\subsection{Position Tracking Error}

\textbf{Motivation.} Because the SO-101 servos accept position setpoints
rather than torques, the gap between commanded and measured position is the
direct, controller-level analogue of a torque error. It captures both
transient failures (the policy commanding more motion than the joint can
deliver under load) and persistent stalls (a joint pinned at its load limit).
Since both signals are recorded synchronously and at the same rate, this is
the cleanest single stability indicator available without extending the
recorder.

Define the per-joint tracking error with a one-step offset to align the
action with the resulting state:
\begin{equation}
    e_i[k] \;=\; a_i[k-1] - q_i[k].
\end{equation}
Per-joint windowed root-mean-square error:
\begin{equation}
    E_i^{\mathrm{RMS}}[k] \;=\; \sqrt{\,\frac{1}{N} \sum_{j \in \mathcal{W}[k]} e_i[j]^2\,}.
\end{equation}
Aggregate norm-error across joints:
\begin{equation}
    E^{\mathrm{RMS}}[k] \;=\; \sqrt{\,\sum_{i=1}^{6} \big(E_i^{\mathrm{RMS}}[k]\big)^2\,}.
\end{equation}

\subsection{Jerk-Based Smoothness}

\textbf{Motivation.} Jerk---the third time-derivative of position---is the
canonical scalar quantification of trajectory roughness, with a long history
in human motor-control literature and in robotics. A policy beginning to lose
stability typically produces jerky motion well before it produces visibly
incorrect motion, providing the earliest-firing indicator in our suite. We
compute jerk on the smoothed measured signal to assess executed motion;
computing it independently on $\mathbf{a}[k]$ (for which no smoothing is
needed) and comparing the two streams further distinguishes
controller-tracking roughness from policy-output roughness.

Discrete third-derivative approximation via a five-point central
finite-difference stencil on the smoothed signal:
\begin{equation}
    \dddot{\tilde{q}}_i[k]
    \;\approx\;
    \frac{\tilde{q}_i[k+2] - 2\,\tilde{q}_i[k+1] + 2\,\tilde{q}_i[k-1] - \tilde{q}_i[k-2]}{2\,T_s^3}.
\end{equation}
Per-joint windowed RMS jerk and its aggregate:
\begin{equation}
    J_i^{\mathrm{RMS}}[k] \;=\; \sqrt{\,\frac{1}{N} \sum_{j \in \mathcal{W}[k]} \dddot{\tilde{q}}_i[j]^2\,},
    \qquad
    J^{\mathrm{RMS}}[k] \;=\; \sqrt{\,\sum_{i=1}^{6} \big(J_i^{\mathrm{RMS}}[k]\big)^2\,}.
\end{equation}

\subsection{Spectral Arc Length (SPARC)}

\textbf{Motivation.} RMS jerk is sensitive to movement scale and duration,
which makes cross-trajectory comparison brittle. SPARC, the negated arc
length of the normalised speed spectrum, is dimensionless and approximately
invariant to amplitude and time-scale changes. It is the most reliable
single-trajectory smoothness statistic in the motor-control literature and
provides a complementary view to RMS jerk: trajectories that look smooth in
RMS but contain narrow-band oscillations are penalised by SPARC.

Define the speed profile
\begin{equation}
    v[k] \;=\; \big\|\, \dot{\tilde{\mathbf{q}}}[k] \,\big\|_2,
\end{equation}
and let $V(\omega)$ denote the discrete-time Fourier transform of $v$ over
the window $\mathcal{W}[k]$, with $\hat{V}(\omega) = V(\omega)/V(0)$ its
normalised magnitude. SPARC is then
\begin{equation}
    \eta_{\mathrm{SPARC}}[k]
    \;=\;
    -\int_{0}^{\omega_c}
    \sqrt{\left(\frac{1}{\omega_c}\right)^{\!2} \!+\! \left(\frac{d\hat{V}(\omega)}{d\omega}\right)^{\!2}}
    \, d\omega,
\end{equation}
with the adaptive cutoff
\begin{equation}
    \omega_c \;=\; \min\!\Big\{\, \omega_c^{\max}, \; \min\!\big\{\omega \,\big|\, \hat{V}(r) < \bar{V}\ \forall\, r > \omega \big\} \,\Big\}.
\end{equation}
At $f_s = 30$~Hz the Nyquist frequency is $15$~Hz, so we use
$\omega_c^{\max} = 12$~Hz and $\bar{V} = 0.05$. Smoother trajectories yield
$\eta_{\mathrm{SPARC}}$ closer to zero from below; rougher motion drives it
more negative.

\subsection{High-Frequency Power Ratio of the Tracking Error}

\textbf{Motivation.} When a position controller cannot follow its setpoint,
the error signal develops oscillatory content well above the bandwidth of the
intended motion---the spectral fingerprint of chatter, limit-cycling, and
incipient instability. On a torque-controlled arm this is conventionally
detected on the commanded torque; lacking torque telemetry, we apply the same
analysis to $e_i[k]$, which carries the same information because closed-loop
oscillation manifests as periodic error around the setpoint.

Let $S_{e_i}(\omega, k)$ denote the Welch power spectral density estimate of
$e_i$ computed on $\mathcal{W}[k]$, with crossover frequency $\omega_h$:
\begin{equation}
    \rho_i^{\mathrm{HF}}[k]
    \;=\;
    \frac{\displaystyle \int_{\omega_h}^{\omega_{\mathrm{Nyq}}} S_{e_i}(\omega, k)\, d\omega}{\displaystyle \int_{0}^{\omega_{\mathrm{Nyq}}} S_{e_i}(\omega, k)\, d\omega}.
\end{equation}
The aggregate is the worst-joint ratio,
\begin{equation}
    \rho^{\mathrm{HF}}[k] \;=\; \max_{i \in \{1,\ldots,6\}} \rho_i^{\mathrm{HF}}[k].
\end{equation}
For SO-101 we use $\omega_h = 5$~Hz and $\omega_{\mathrm{Nyq}} = 15$~Hz.
Sustained $\rho^{\mathrm{HF}} > 0.2$ is a strong chatter indicator on
arm-scale dynamics.

\subsection{Stall Indicator and Joint-Limit Clipping}

\textbf{Motivation.} On a torque-instrumented arm we would threshold the
commanded torque against the motor's torque limit. Without that telemetry we
detect saturation behaviourally: when the policy commands meaningful motion
but the joint fails to move, the motor is effectively at its load limit. We
combine this with a separate detector for action saturation against the
calibrated joint range.

Let $\delta_a^{(i)}, \delta_q^{(i)}$ be per-joint command and motion
thresholds, calibrated as the 95th percentile of the corresponding magnitudes
on a corpus of successful teacher rollouts. Over a short look-back of $M$
samples, the instantaneous stall indicator is
\begin{equation}
    \sigma_i^{\mathrm{stall}}[k] \;=\; \mathbbm{1}\!\left[\; |a_i[k] - a_i[k-M]| > \delta_a^{(i)} \;\wedge\; |q_i[k] - q_i[k-M]| < \delta_q^{(i)} \;\right].
\end{equation}
With calibrated joint range $[a_i^{\min}, a_i^{\max}]$ and tolerance
$\epsilon_i = 0.02 \cdot (a_i^{\max} - a_i^{\min})$, the clipping indicator is
\begin{equation}
    \sigma_i^{\mathrm{clip}}[k] \;=\; \mathbbm{1}\!\left[\; a_i[k] \le a_i^{\min} + \epsilon_i \;\vee\; a_i[k] \ge a_i^{\max} - \epsilon_i \;\right].
\end{equation}
The combined per-joint flag and its windowed aggregate are
\begin{align}
    \sigma_i[k] &\;=\; \sigma_i^{\mathrm{stall}}[k] \;\vee\; \sigma_i^{\mathrm{clip}}[k], \\
    \bar{\sigma}[k] &\;=\; \frac{1}{N} \sum_{j \in \mathcal{W}[k]} \mathbbm{1}\!\left[\, \max_{i} \sigma_i[j] = 1 \,\right].
\end{align}
We use $M = 5$ ($\approx 167$~ms at 30~Hz), short enough to react to abrupt
stalls but long enough to ignore single-frame quantisation jitter.

\subsection{Combined Intervention Rule}

We assemble the five sign-aligned signals (larger means worse) into a single
monitor vector,
\begin{equation}
    \mathbf{m}[k] \;=\; \big[\, E^{\mathrm{RMS}}[k], \; J^{\mathrm{RMS}}[k], \; -\eta_{\mathrm{SPARC}}[k], \; \rho^{\mathrm{HF}}[k], \; \bar{\sigma}[k] \,\big]^\top \in \mathbb{R}^5.
\end{equation}
A threshold vector $\bm{\theta} \in \mathbb{R}^5$ is calibrated as the 95th
percentile of each component over a held-out corpus of successful teacher
rollouts. The any-of intervention rule fires when any single channel
exceeds its threshold,
\begin{equation}
    \mathcal{I}^{\mathrm{any}}[k]
    \;=\;
    \mathbbm{1}\!\left[\, \exists\, \ell \in \{1,\ldots,5\} : m_\ell[k] > \theta_\ell \,\right],
\end{equation}
while the weighted-sum rule, with non-negative weights
$\mathbf{w} \in \mathbb{R}_{\ge 0}^{5}$ and aggregate threshold
$\theta_{\Sigma}$, integrates evidence across channels:
\begin{equation}
    \mathcal{I}^{\Sigma}[k]
    \;=\;
    \mathbbm{1}\!\left[\, \sum_{\ell=1}^{5} w_\ell \, \frac{m_\ell[k]}{\theta_\ell} > \theta_{\Sigma} \,\right].
\end{equation}
The any-of rule fires earlier and is easier to attribute when debugging; the
weighted-sum rule is more robust to single-channel noise. To suppress
single-frame artefacts, we additionally require the chosen rule to remain
active for $P$ consecutive samples before triggering an intervention,
\begin{equation}
    \mathcal{I}^{\star}[k] \;=\; \prod_{p = 0}^{P-1} \mathcal{I}[k - p],
\end{equation}
with $P = 3$ at 30~Hz ($\approx 100$~ms persistence).

\paragraph{Optional extension.} If \texttt{Present\_Load} or
\texttt{Present\_Current} is added to the recorded observation features by
extending the LeRobot \texttt{Robot} interface, the stall indicator
$\sigma_i^{\mathrm{stall}}$ can be replaced by the conventional torque
saturation rate, and the high-frequency power ratio can be recomputed on the
load signal directly, recovering the full classical signal set without
modifying the rest of the pipeline.
, or paste directly.
%
%   Required packages in your main preamble:
%       \usepackage{amsmath, amssymb, amsfonts}
%       \usepackage{bbm}    % for \mathbbm{1} indicator function
%       \usepackage{bm}     % for \bm{} bold math
%
%   If \mathbbm is unavailable in your TeX distribution, substitute:
%       \usepackage{dsfont}     and use \mathds{1}
%   or fall back to \mathbf{1} (no extra package).
% =============================================================================

\section{Methods: Online Stability Monitoring}

We monitor the rollouts of a vision--language--action (VLA) student policy on
the 6-DOF SO-101 arm and intervene when classical control-theoretic
indicators of trajectory instability exceed thresholds calibrated on a teacher
policy's successful rollouts. All indicators are derived solely from the data
that the LeRobot \texttt{record} pipeline persists by default, so the monitor
imposes no additional logging requirements and runs at a small fraction of the
cost of a second VLA inference pass.

\subsection{Notation and Data}

The SO-101 has six Feetech STS3215 serial-bus servos, indexed
$i \in \{1,\ldots,6\}$ in the order
\textit{shoulder\_pan}, \textit{shoulder\_lift}, \textit{elbow\_flex},
\textit{wrist\_flex}, \textit{wrist\_roll}, \textit{gripper}.
The default \texttt{lerobot-record} schema persists, at sampling rate
$f_s = 30$~Hz with period $T_s = 1/f_s$, two synchronized vectors per frame:
\begin{align}
    \mathbf{a}[k] &= \big[a_1[k], \ldots, a_6[k]\big]^\top \in \mathbb{R}^6 && \text{(commanded joint positions; \texttt{action})} \\
    \mathbf{q}[k] &= \big[q_1[k], \ldots, q_6[k]\big]^\top \in \mathbb{R}^6 && \text{(measured joint positions; \texttt{observation.state})}
\end{align}
The motors expose load, current, velocity, and temperature feedback over the
serial bus, but these channels are not logged by the default recorder; the
metrics below therefore use $\mathbf{a}$ and $\mathbf{q}$ exclusively.

For real-time monitoring we compute all metrics over a sliding window of
length $N$ samples ending at the current step $k$,
\begin{equation}
    \mathcal{W}[k] \;=\; \{k - N + 1, \ldots, k\}, \qquad \Delta \;=\; N T_s.
\end{equation}
The 12-bit magnetic encoder on the STS3215 quantises position at $\approx
0.088^\circ$ and the gearbox introduces up to $0.5^\circ$ of backlash, so
naive numerical differentiation of $\mathbf{q}$ is dominated by quantisation
noise. We therefore apply a Savitzky--Golay smoother (window 7, polynomial
order 3) to $\mathbf{q}[k]$ before differentiation, denoting the smoothed
signal $\tilde{\mathbf{q}}[k]$.

\subsection{Position Tracking Error}

\textbf{Motivation.} Because the SO-101 servos accept position setpoints
rather than torques, the gap between commanded and measured position is the
direct, controller-level analogue of a torque error. It captures both
transient failures (the policy commanding more motion than the joint can
deliver under load) and persistent stalls (a joint pinned at its load limit).
Since both signals are recorded synchronously and at the same rate, this is
the cleanest single stability indicator available without extending the
recorder.

Define the per-joint tracking error with a one-step offset to align the
action with the resulting state:
\begin{equation}
    e_i[k] \;=\; a_i[k-1] - q_i[k].
\end{equation}
Per-joint windowed root-mean-square error:
\begin{equation}
    E_i^{\mathrm{RMS}}[k] \;=\; \sqrt{\,\frac{1}{N} \sum_{j \in \mathcal{W}[k]} e_i[j]^2\,}.
\end{equation}
Aggregate norm-error across joints:
\begin{equation}
    E^{\mathrm{RMS}}[k] \;=\; \sqrt{\,\sum_{i=1}^{6} \big(E_i^{\mathrm{RMS}}[k]\big)^2\,}.
\end{equation}

\subsection{Jerk-Based Smoothness}

\textbf{Motivation.} Jerk---the third time-derivative of position---is the
canonical scalar quantification of trajectory roughness, with a long history
in human motor-control literature and in robotics. A policy beginning to lose
stability typically produces jerky motion well before it produces visibly
incorrect motion, providing the earliest-firing indicator in our suite. We
compute jerk on the smoothed measured signal to assess executed motion;
computing it independently on $\mathbf{a}[k]$ (for which no smoothing is
needed) and comparing the two streams further distinguishes
controller-tracking roughness from policy-output roughness.

Discrete third-derivative approximation via a five-point central
finite-difference stencil on the smoothed signal:
\begin{equation}
    \dddot{\tilde{q}}_i[k]
    \;\approx\;
    \frac{\tilde{q}_i[k+2] - 2\,\tilde{q}_i[k+1] + 2\,\tilde{q}_i[k-1] - \tilde{q}_i[k-2]}{2\,T_s^3}.
\end{equation}
Per-joint windowed RMS jerk and its aggregate:
\begin{equation}
    J_i^{\mathrm{RMS}}[k] \;=\; \sqrt{\,\frac{1}{N} \sum_{j \in \mathcal{W}[k]} \dddot{\tilde{q}}_i[j]^2\,},
    \qquad
    J^{\mathrm{RMS}}[k] \;=\; \sqrt{\,\sum_{i=1}^{6} \big(J_i^{\mathrm{RMS}}[k]\big)^2\,}.
\end{equation}

\subsection{Spectral Arc Length (SPARC)}

\textbf{Motivation.} RMS jerk is sensitive to movement scale and duration,
which makes cross-trajectory comparison brittle. SPARC, the negated arc
length of the normalised speed spectrum, is dimensionless and approximately
invariant to amplitude and time-scale changes. It is the most reliable
single-trajectory smoothness statistic in the motor-control literature and
provides a complementary view to RMS jerk: trajectories that look smooth in
RMS but contain narrow-band oscillations are penalised by SPARC.

Define the speed profile
\begin{equation}
    v[k] \;=\; \big\|\, \dot{\tilde{\mathbf{q}}}[k] \,\big\|_2,
\end{equation}
and let $V(\omega)$ denote the discrete-time Fourier transform of $v$ over
the window $\mathcal{W}[k]$, with $\hat{V}(\omega) = V(\omega)/V(0)$ its
normalised magnitude. SPARC is then
\begin{equation}
    \eta_{\mathrm{SPARC}}[k]
    \;=\;
    -\int_{0}^{\omega_c}
    \sqrt{\left(\frac{1}{\omega_c}\right)^{\!2} \!+\! \left(\frac{d\hat{V}(\omega)}{d\omega}\right)^{\!2}}
    \, d\omega,
\end{equation}
with the adaptive cutoff
\begin{equation}
    \omega_c \;=\; \min\!\Big\{\, \omega_c^{\max}, \; \min\!\big\{\omega \,\big|\, \hat{V}(r) < \bar{V}\ \forall\, r > \omega \big\} \,\Big\}.
\end{equation}
At $f_s = 30$~Hz the Nyquist frequency is $15$~Hz, so we use
$\omega_c^{\max} = 12$~Hz and $\bar{V} = 0.05$. Smoother trajectories yield
$\eta_{\mathrm{SPARC}}$ closer to zero from below; rougher motion drives it
more negative.

\subsection{High-Frequency Power Ratio of the Tracking Error}

\textbf{Motivation.} When a position controller cannot follow its setpoint,
the error signal develops oscillatory content well above the bandwidth of the
intended motion---the spectral fingerprint of chatter, limit-cycling, and
incipient instability. On a torque-controlled arm this is conventionally
detected on the commanded torque; lacking torque telemetry, we apply the same
analysis to $e_i[k]$, which carries the same information because closed-loop
oscillation manifests as periodic error around the setpoint.

Let $S_{e_i}(\omega, k)$ denote the Welch power spectral density estimate of
$e_i$ computed on $\mathcal{W}[k]$, with crossover frequency $\omega_h$:
\begin{equation}
    \rho_i^{\mathrm{HF}}[k]
    \;=\;
    \frac{\displaystyle \int_{\omega_h}^{\omega_{\mathrm{Nyq}}} S_{e_i}(\omega, k)\, d\omega}{\displaystyle \int_{0}^{\omega_{\mathrm{Nyq}}} S_{e_i}(\omega, k)\, d\omega}.
\end{equation}
The aggregate is the worst-joint ratio,
\begin{equation}
    \rho^{\mathrm{HF}}[k] \;=\; \max_{i \in \{1,\ldots,6\}} \rho_i^{\mathrm{HF}}[k].
\end{equation}
For SO-101 we use $\omega_h = 5$~Hz and $\omega_{\mathrm{Nyq}} = 15$~Hz.
Sustained $\rho^{\mathrm{HF}} > 0.2$ is a strong chatter indicator on
arm-scale dynamics.

\subsection{Stall Indicator and Joint-Limit Clipping}

\textbf{Motivation.} On a torque-instrumented arm we would threshold the
commanded torque against the motor's torque limit. Without that telemetry we
detect saturation behaviourally: when the policy commands meaningful motion
but the joint fails to move, the motor is effectively at its load limit. We
combine this with a separate detector for action saturation against the
calibrated joint range.

Let $\delta_a^{(i)}, \delta_q^{(i)}$ be per-joint command and motion
thresholds, calibrated as the 95th percentile of the corresponding magnitudes
on a corpus of successful teacher rollouts. Over a short look-back of $M$
samples, the instantaneous stall indicator is
\begin{equation}
    \sigma_i^{\mathrm{stall}}[k] \;=\; \mathbbm{1}\!\left[\; |a_i[k] - a_i[k-M]| > \delta_a^{(i)} \;\wedge\; |q_i[k] - q_i[k-M]| < \delta_q^{(i)} \;\right].
\end{equation}
With calibrated joint range $[a_i^{\min}, a_i^{\max}]$ and tolerance
$\epsilon_i = 0.02 \cdot (a_i^{\max} - a_i^{\min})$, the clipping indicator is
\begin{equation}
    \sigma_i^{\mathrm{clip}}[k] \;=\; \mathbbm{1}\!\left[\; a_i[k] \le a_i^{\min} + \epsilon_i \;\vee\; a_i[k] \ge a_i^{\max} - \epsilon_i \;\right].
\end{equation}
The combined per-joint flag and its windowed aggregate are
\begin{align}
    \sigma_i[k] &\;=\; \sigma_i^{\mathrm{stall}}[k] \;\vee\; \sigma_i^{\mathrm{clip}}[k], \\
    \bar{\sigma}[k] &\;=\; \frac{1}{N} \sum_{j \in \mathcal{W}[k]} \mathbbm{1}\!\left[\, \max_{i} \sigma_i[j] = 1 \,\right].
\end{align}
We use $M = 5$ ($\approx 167$~ms at 30~Hz), short enough to react to abrupt
stalls but long enough to ignore single-frame quantisation jitter.

\subsection{Combined Intervention Rule}

We assemble the five sign-aligned signals (larger means worse) into a single
monitor vector,
\begin{equation}
    \mathbf{m}[k] \;=\; \big[\, E^{\mathrm{RMS}}[k], \; J^{\mathrm{RMS}}[k], \; -\eta_{\mathrm{SPARC}}[k], \; \rho^{\mathrm{HF}}[k], \; \bar{\sigma}[k] \,\big]^\top \in \mathbb{R}^5.
\end{equation}
A threshold vector $\bm{\theta} \in \mathbb{R}^5$ is calibrated as the 95th
percentile of each component over a held-out corpus of successful teacher
rollouts. The any-of intervention rule fires when any single channel
exceeds its threshold,
\begin{equation}
    \mathcal{I}^{\mathrm{any}}[k]
    \;=\;
    \mathbbm{1}\!\left[\, \exists\, \ell \in \{1,\ldots,5\} : m_\ell[k] > \theta_\ell \,\right],
\end{equation}
while the weighted-sum rule, with non-negative weights
$\mathbf{w} \in \mathbb{R}_{\ge 0}^{5}$ and aggregate threshold
$\theta_{\Sigma}$, integrates evidence across channels:
\begin{equation}
    \mathcal{I}^{\Sigma}[k]
    \;=\;
    \mathbbm{1}\!\left[\, \sum_{\ell=1}^{5} w_\ell \, \frac{m_\ell[k]}{\theta_\ell} > \theta_{\Sigma} \,\right].
\end{equation}
The any-of rule fires earlier and is easier to attribute when debugging; the
weighted-sum rule is more robust to single-channel noise. To suppress
single-frame artefacts, we additionally require the chosen rule to remain
active for $P$ consecutive samples before triggering an intervention,
\begin{equation}
    \mathcal{I}^{\star}[k] \;=\; \prod_{p = 0}^{P-1} \mathcal{I}[k - p],
\end{equation}
with $P = 3$ at 30~Hz ($\approx 100$~ms persistence).

\paragraph{Optional extension.} If \texttt{Present\_Load} or
\texttt{Present\_Current} is added to the recorded observation features by
extending the LeRobot \texttt{Robot} interface, the stall indicator
$\sigma_i^{\mathrm{stall}}$ can be replaced by the conventional torque
saturation rate, and the high-frequency power ratio can be recomputed on the
load signal directly, recovering the full classical signal set without
modifying the rest of the pipeline.
